% GENERATED FILE. Edit canonical lessons, then run npm run generate:books.
% canonical-source-hash: fnv1a64:e08e1a96edaa5fe4
% canonical-lessons: IT-C513-gareggiare, IT-C513-cronometro, IT-C513-maratona, IT-C513-atleta, IT-C513-ansimare

\chapter{La gara}
\label{ch:it-la-gara}
\hlchaptermodality{513}

\begin{chapteropening}
\textbf{By the end of this chapter:} \emph{I can describe a race: timing, distance and being out of breath.}

Complete the last lesson of chapter 513: I can describe a race: timing, distance and being out of breath.
\end{chapteropening}

\section[gareggiare]{gareggiare — to compete, to race}

\label{lesson:IT-C513-gareggiare}



\begin{quote}
Before the new one: say the Italian for medal, then the Italian for podium.
\end{quote}

\subsection*{You'll want to know: gareggiare}
\textbf{gareggiare} — \textquotedblleft{}to compete, to race\textquotedblright{}.

\begin{quote}
Gareggio nei cento metri.
\end{quote}

\begin{quote}
\emph{I compete in the hundred metres.}
\end{quote}

Built on \emph{la gara} \textquotedblleft{}competition\textquotedblright{}.

\begin{cousinweb}
\emph{gareggiare} is built on \emph{gara} \textquotedblleft{}a competition\textquotedblright{}, whose origin is not settled; an Arabic source is often suggested.
\end{cousinweb}

\begin{grammarlens}[title={What you've built}]
One more word for telling, explaining and describing at B1.
\end{grammarlens}

\subsection*{Your turn}
\begin{itemize}
\raggedright
  \item \emph{Say it:} \emph{gareggiare}
  \item \emph{Say it:} \emph{gareggiare}, once more
  \item \emph{Say it:} a sentence of your own with \emph{gareggiare}
  \item \emph{Recall:} say the Italian for medal, then the Italian for podium, then say \emph{gareggiare} again
\end{itemize}

\begin{tcolorbox}[breakable,colback=teal!4,colframe=teal!35!black,title={Before you move on}]
What does gareggiare mean? (\textquotedblleft{}To compete, to race\textquotedblright{}.) Say it once more.
\end{tcolorbox}
\section[il cronometro]{il cronometro — stopwatch}

\label{lesson:IT-C513-cronometro}



\begin{quote}
Before the new one: say the Italian for podium, then the Italian for to compete.
\end{quote}

\subsection*{You'll want to know: il cronometro}
\textbf{il cronometro} — \textquotedblleft{}stopwatch\textquotedblright{}.

\begin{quote}
L'allenatore ferma il cronometro all'arrivo.
\end{quote}

\begin{quote}
\emph{The coach stops the stopwatch at the finish.}
\end{quote}

\emph{Una gara a cronometro} — a time trial.

\begin{cousinweb}
\emph{cronometro} is a modern learned word built from Greek \emph{khrónos} \textquotedblleft{}time\textquotedblright{} and \emph{métron} \textquotedblleft{}measure\textquotedblright{} (EtymDB).
\end{cousinweb}

\begin{grammarlens}[title={What you've built}]
One more word for telling, explaining and describing at B1.
\end{grammarlens}

\subsection*{Your turn}
\begin{itemize}
\raggedright
  \item \emph{Say it:} \emph{il cronometro}
  \item \emph{Say it:} \emph{il cronometro}, once more
  \item \emph{Say it:} a sentence of your own with \emph{il cronometro}
  \item \emph{Recall:} say the Italian for podium, then the Italian for to compete, then say \emph{il cronometro} again
\end{itemize}

\begin{tcolorbox}[breakable,colback=teal!4,colframe=teal!35!black,title={Before you move on}]
What does il cronometro mean? (\textquotedblleft{}Stopwatch\textquotedblright{}.) Say it once more.
\end{tcolorbox}
\section[la maratona]{la maratona — marathon}

\label{lesson:IT-C513-maratona}



\begin{quote}
Before the new one: say the Italian for to compete, then the Italian for stopwatch.
\end{quote}

\subsection*{You'll want to know: la maratona}
\textbf{la maratona} — \textquotedblleft{}marathon\textquotedblright{}.

\begin{quote}
Ho corso la mia prima maratona a quarant'anni.
\end{quote}

\begin{quote}
\emph{I ran my first marathon at forty.}
\end{quote}

Also a long session: \emph{una maratona di film}.

\begin{cousinweb}
\emph{maratona} comes from Greek, from the name of a town, which is built on \emph{márathon} \textquotedblleft{}fennel\textquotedblright{} (EtymDB): the story goes that a runner carried news of a victory from there.
\end{cousinweb}

\begin{grammarlens}[title={What you've built}]
One more word for telling, explaining and describing at B1.
\end{grammarlens}

\subsection*{Your turn}
\begin{itemize}
\raggedright
  \item \emph{Say it:} \emph{la maratona}
  \item \emph{Say it:} \emph{la maratona}, once more
  \item \emph{Say it:} a sentence of your own with \emph{la maratona}
  \item \emph{Recall:} say the Italian for to compete, then the Italian for stopwatch, then say \emph{la maratona} again
\end{itemize}

\begin{tcolorbox}[breakable,colback=teal!4,colframe=teal!35!black,title={Before you move on}]
What does la maratona mean? (\textquotedblleft{}Marathon\textquotedblright{}.) Say it once more.
\end{tcolorbox}
\section[l'atleta]{l'atleta — athlete}

\label{lesson:IT-C513-atleta}



\begin{quote}
Before the new one: say the Italian for stopwatch, then the Italian for marathon.
\end{quote}

\subsection*{You'll want to know: l'atleta}
\textbf{l'atleta} — \textquotedblleft{}athlete\textquotedblright{}.

\begin{quote}
Gli atleti si preparano per la gara.
\end{quote}

\begin{quote}
\emph{The athletes are getting ready for the race.}
\end{quote}

The same form for men and women: \emph{l'atleta}; plural \emph{gli atleti}, \emph{le atlete}.

\begin{cousinweb}
\emph{atleta} is a learned borrowing from Latin \emph{āthlēta} (EtymDB), from Greek \emph{athlēt\'{\={e}}s} \textquotedblleft{}one who competes for a prize\textquotedblright{}, from \emph{âthlos} \textquotedblleft{}a contest\textquotedblright{}.
\end{cousinweb}

\begin{grammarlens}[title={What you've built}]
One more word for telling, explaining and describing at B1.
\end{grammarlens}

\subsection*{Your turn}
\begin{itemize}
\raggedright
  \item \emph{Say it:} \emph{l'atleta}
  \item \emph{Say it:} \emph{l'atleta}, once more
  \item \emph{Say it:} a sentence of your own with \emph{l'atleta}
  \item \emph{Recall:} say the Italian for stopwatch, then the Italian for marathon, then say \emph{l'atleta} again
\end{itemize}

\begin{tcolorbox}[breakable,colback=teal!4,colframe=teal!35!black,title={Before you move on}]
What does l'atleta mean? (\textquotedblleft{}Athlete\textquotedblright{}.) Say it once more.
\end{tcolorbox}
\section[ansimare]{ansimare — to pant, to gasp}

\label{lesson:IT-C513-ansimare}



\begin{quote}
Before the new one: say the Italian for marathon, then the Italian for athlete.
\end{quote}

\subsection*{You'll want to know: ansimare}
\textbf{ansimare} — \textquotedblleft{}to pant, to gasp\textquotedblright{}.

\begin{quote}
All'arrivo ansimavo e non riuscivo a parlare.
\end{quote}

\begin{quote}
\emph{At the finish I was panting and could not speak.}
\end{quote}

Breathing hard after an effort.

\begin{cousinweb}
\emph{ansimare} is usually traced to Latin \emph{anhēlāre} \textquotedblleft{}to pant\textquotedblright{} (Lewis \& Short), through a Vulgar Latin form; the details are not settled.
\end{cousinweb}

\begin{grammarlens}[title={What you've built}]
One more word for telling, explaining and describing at B1.
\end{grammarlens}

\subsection*{Your turn}
\begin{itemize}
\raggedright
  \item \emph{Say it:} \emph{ansimare}
  \item \emph{Say it:} \emph{ansimare}, once more
  \item \emph{Say it:} a sentence of your own with \emph{ansimare}
  \item \emph{Recall:} say the Italian for marathon, then the Italian for athlete, then say \emph{ansimare} again
\end{itemize}

\begin{tcolorbox}[breakable,colback=teal!4,colframe=teal!35!black,title={Before you move on}]
What does ansimare mean? (\textquotedblleft{}To pant, to gasp\textquotedblright{}.) Say it once more.
\end{tcolorbox}
